\chapter{Modèle mathématique et difficulté du problème}
\section{Une solution sous forme d'anneau et d'affectations}
Soit $V=\{1,\ldots,n\}$ l'ensemble des points et $r\in V$ le sommet imposé.
Pour des coordonnées $(x_i,y_i)$, le projet définit
\begin{equation}
d_{ij}=\sqrt{(x_i-x_j)^2+(y_i-y_j)^2}.
\end{equation}
Ces distances sont symétriques, positives ou nulles, et non arrondies. Le graphe
est complet : une arête peut relier n'importe quelle paire de points.

Une solution est un triplet $(S,\pi,a)$. L'ensemble $S\subseteq V$ contient les
$p$ stations et vérifie $r\in S$. Le vecteur $\pi=(\pi_1,\ldots,\pi_p)$ est une
permutation de $S$, qui définit les arêtes du cycle en posant $\pi_{p+1}=\pi_1$.
L'affectation $a(i)\in S$ associe chaque point à une station, avec $a(i)=i$
pour les stations.

Le cycle est \emph{simple au sens combinatoire} : aucune station n'est répétée
dans la permutation. Cette propriété n'interdit pas, à elle seule, un croisement
géométrique de deux segments. Le recours à 2-opt améliore le parcours euclidien,
sans ajouter explicitement une contrainte géométrique d'absence de croisement.

\section{Fonction objectif et affectation optimale à stations fixées}
Les deux composantes sont
\begin{equation}
L(S,\pi)=\sum_{k=1}^{p}d_{\pi_k,\pi_{k+1}},\qquad
A(S,a)=\sum_{i\in V}d_{i,a(i)}.
\end{equation}
Le problème minimise
\begin{equation}\label{eq:objectif}
C(S,\pi,a)=\alpha L(S,\pi)+A(S,a).
\end{equation}
Les stations contribuent zéro à $A$, puisqu'elles s'affectent à elles-mêmes.
La longueur $L$ est comptée une seule fois sur le cycle fermé. Il ne s'agit pas
de la somme des trajets en transport de tous les voyageurs : aucun couple
origine-destination n'est modélisé.

Lorsque $S$ est fixé, les affectations sont indépendantes les unes des autres.
Il n'existe ni capacité de station ni quota d'usagers. On peut donc choisir
\begin{equation}
a(i)\in\mathop{\mathrm{arg\,min}}_{s\in S}d_{is},\qquad
A(S)=\sum_{i\in V}\min_{s\in S}d_{is}.
\end{equation}
Toute affectation plus lointaine peut être remplacée par une plus proche sans
modifier l'anneau et sans augmenter le coût. Les heuristiques exploitent cette
propriété à chaque échange de stations. En cas d'égalité de distance, elles
choisissent le plus petit indice pour rendre la décision déterministe.

\section{Deux problèmes classiques intégrés}
La sélection des stations rappelle le problème du $p$-médian, qui minimise la
somme des distances d'accès. Une fois les stations fixées, la recherche du
meilleur ordre de visite est un problème du voyageur de commerce sur $S$.
L'intégration vient du fait que choisir $S$ en ne regardant que les accès peut
produire un anneau coûteux. La meilleure sélection dépend donc de $\alpha$.

Pour $\alpha=0$, seul le coût d'accès intervient. Le problème devient un
$p$-médian avec un centre imposé ; le cycle reste requis pour la représentation,
mais son ordre n'influe plus sur l'objectif. Pour $p=n$, toutes les affectations
ont un coût nul. Si $\alpha>0$, il ne reste que le choix du cycle.

Une augmentation de $\alpha$ renchérit toute solution fixée dont l'anneau est
non nul, mais peut aussi conduire à sélectionner d'autres stations. De même,
augmenter $p$ ne garantit pas une diminution du coût total : une meilleure
couverture et une ligne plus longue peuvent se compenser. Les graphiques selon
$p$ doivent ainsi être lus comme des compromis, et non comme une courbe de
performance qui devrait forcément décroître.

\section{Justification de la NP-difficulté}
Considérons une instance du voyageur de commerce euclidien symétrique à $n$
points, avec $n\geq3$. Construisons une instance Ring-Star ayant les mêmes points,
les mêmes distances, $p=n$, $\alpha=1$ et un sommet imposé quelconque.
Cette transformation demande un temps polynomial.

Comme exactement $n$ stations sont requises, tous les points sont sélectionnés.
Le coût des affectations est nul et le coût Ring-Star est exactement la longueur
du cycle hamiltonien. Le sommet imposé ne retire aucune solution, puisque tout
cycle hamiltonien visite tous les sommets. Une résolution optimale du Ring-Star
donnerait donc une résolution optimale du voyageur de commerce euclidien.
Ce dernier étant NP-difficile, le problème Ring-Star général l'est également.

Cette réduction porte sur un cas particulier autorisé, avec $\alpha=1$.
Elle ne dit pas que chaque instance est difficile ni que le cas dégénéré
$\alpha=0$, $p=n$ l'est. Elle justifie le recours aux heuristiques et la nécessité
de limiter certaines résolutions exactes. Il suffit ici d'établir la
NP-difficulté ; aucune affirmation supplémentaire sur l'appartenance à NP avec
des distances euclidiennes exactes n'est nécessaire.
\bilan{Le coût combine deux décisions liées. Les affectations sont simples à
optimiser lorsque les stations sont fixées, mais la sélection des stations et
l'ordre de leur cycle rendent le problème combinatoire.}
